% 01-intro.tex — §1 引言
\section{引言}

区块链信息系统的执行层本质上是一台多副本共识状态机：所有参与共识的节点必须对
同一笔交易产生按位一致的执行结果。以太坊虚拟机（EVM）是当前最广泛部署的智能
合约执行环境，其即时编译器（JIT）已在多种主流客户端中逐步取代解释器，成为
共识关键路径上的核心组件。JIT 编译器中的窥孔重写规则（peephole，代码生成路径上的局部模式匹配）直接作用于
机器码生成阶段，其正确性直接关系到链上执行的确定性保证。

然而，若某条窥孔重写规则的等价性判断有误，且只在部分客户端上线，就可能导致
节点间产生不一致的执行结果，进而引发链分叉——这是区块链信息系统中代价最高的
故障模式之一。现有形式化窥孔验证工作已对 EVM 优化做出重要贡献：例如
Albert 等人~\cite{verified_evm_cav23}在字节码块级提出了经 Coq 形式化验证的
优化规则集合。但该工作关注的是 EVM 字节码层的抽象，与 JIT 编译器在代码生成
阶段产出的机器级模式——尤其是 256 位字长映射到 64 位硬件后形成的多字进位链
（由 \texttt{adc}/\texttt{sbb} 指令构成）——位于相邻但不同的抽象层次，
对后者的覆盖不足。

这一空白的根源在于中间表示对进位位的建模。在通用编译器中间表示（如 LLVM IR）中，
基础算术运算仅对 $N$ 位整数定义；进位位需通过返回 \{结果, 进位位\} 多返回值结构的
内置函数（如 \texttt{@llvm.uadd.with.overflow}）间接获取——进位位不是一等公民。
因此，依赖进位位的多字规则在自动等价验证工具（如 Alive~\cite{alive}）中只能借助
间接结构表达；此时的瓶颈不在 SMT 求解器~\cite{demoura2008z3}的容量，而在 IR
的模式匹配能力。
本文在 DTVM 的 dMIR（确定性中间表示，为共识确定性放弃与通用 IR 的直接互操作，
详见第~\ref{sec:two-level}~节）中把带进位加减法
\texttt{adc}/\texttt{sbb} 作为一等位向量算子，
从而把这类以多字进位链为核心的规则纳入 SMT 可自动验证的范围。

\paragraph{贡献。}本文的具体贡献如下：
\begin{enumerate}
  \item 构建了面向确定性 EVM JIT 的两级窥孔重写系统：在 dMIR 中把
    \texttt{adc}/\texttt{sbb} 作为一等位向量算子，共得到 \textbf{70} 条
    dMIR 规则与 \textbf{13} 条 x86 代码生成层规则；dMIR 层由 Z3 位向量等价性
    校验~\cite{demoura2008z3}，x86 层由对照测试套件覆盖。
  \item 揭示与已有形式化 EVM 优化工作的层次互补：与
    Albert 等人~\cite{verified_evm_cav23}的 EVM 字节码块级优化规则集合相比，
    规则重合度仅 \textbf{8.4\%}（7/83 条），其余 76 条分布在字节码层无法
    触及的抽象层次。
  \item 在 27 个标准 EVM 基准上验证性能与正确性：几何均值加速
    \textbf{+5.93\%}（95\% 置信区间 [+3.53\%, +8.78\%]），峰值
    \textbf{+24.7\%}；3161 条测试全部通过且开启/关闭规则集两侧结果一致。
\end{enumerate}
